\begin{lemma}[Higher symmetric terms on the upper critical line]
\label{mod:cases-wildodd-highervanishing}
Let $K/F$ be a wildly ramified cyclic extension of odd prime degree
\[
 p=[K:F]=\operatorname{char}k_F,
\]
with lower break $t>0$, and put $T=t+1$.  Suppose that the residue degree
is one and retain the integral generating uniformizer used in the exact
trace-ideal and norm-pullback conductor theorems.  Write
\[
 m_F(\chi_F)=2d+1>1,\qquad b=2d,
\]
and define the signed upper shift
\begin{equation}\label{eq:lean-wildodd-upper-shift}
 a=\texttt{wildOddUpperShift}(t,d)=(t:\mathbf Z)-2(d:\mathbf Z)=t-b.
\end{equation}
The exact last nontrivial upper layer is
\begin{equation}\label{eq:lean-wildodd-upper-depth}
 \texttt{wildOddUpperDepth}(p,t,d)=
 \begin{cases}
  2d,&2d<t,\\
  t,&2d=t,\\
  t+p(2d-t),&t<2d.
 \end{cases}
\end{equation}
The predicate $\texttt{WildOddUpperDepthFormula}(p,t,d,d_K)$ keeps these
three rows as separate implications:
\[
 2d<t\Longrightarrow d_K=d,\qquad
 2d=t\Longrightarrow d_K=d,
\]
and
\[
 t<2d\Longrightarrow
 2d_K=t+p(2d-t)
\]
as an equality in $\mathbf Z$.  Its theorem \texttt{eq\_piecewise}
recovers
\[
 2d_K=\texttt{wildOddUpperDepth}(p,t,d).
\]

For $u,x\in K$ and $j\in\mathbf N$, put
\begin{align}
 C_j(u;x)
  &=\texttt{wildOddUpperCorrectionTerm}(F,K,u,x,j)\notag\\
  &=s_j(ux)-N_{K/F}(u)s_j(x),
 \label{eq:lean-wildodd-correction-term}\\
 C(u;x)
  &=\texttt{wildOddUpperCorrection}(F,K,p,u,x)\notag\\
  &=\sum_{j\in\operatorname{Icc}(2,p)}C_j(u;x),
 \label{eq:lean-wildodd-correction}\\
 Q(u;x)
  &=\texttt{wildOddUpperQuadraticCorrection}(F,K,u,x)=C_2(u;x).
 \label{eq:lean-wildodd-quadratic}
\end{align}
Thus the scalar denoted $n$ in the manuscript is definitionally
$N_{K/F}(u)$; it is never an unrelated parameter.

The coefficient function is extended by literal zero outside degrees
$2,\ldots,p$:
\begin{equation}\label{eq:lean-wildodd-coefficient}
 \texttt{wildOddUpperCorrectionCoefficient}(F,K,p,u,x,j)=
 \begin{cases}
  C_j(u;x),&2\le j\le p,\\
  0,&\text{otherwise}.
 \end{cases}
\end{equation}
Consequently the degree-$0$ and degree-$1$ coefficients are zero, while
the degree-$2$ coefficient is exactly $Q(u;x)$.  Lean records these as
three separate theorems.  In particular, absence of constant and affine
terms is not inferred from an informal description of the sum.

If $p=[K:F]$, the terminal coefficient cancels by the exact equality
\begin{equation}\label{eq:lean-wildodd-endpoint-cancellation}
 C_p(u;x)=s_p(ux)-N_{K/F}(u)s_p(x)
 =N_{K/F}(ux)-N_{K/F}(u)N_{K/F}(x)=0.
\end{equation}
This is \texttt{wildOddUpperCorrectionTerm\_degree}; it is not replaced
by lattice membership.  When $p=3$, the range $3\le j<p$ is empty, as
recorded separately by
\texttt{wildOdd\_intermediate\_range\_empty\_degree\_three}.  For
$p\ge3$, the exact finite-sum decomposition is
\begin{equation}\label{eq:lean-wildodd-correction-decomposition}
 C(u;x)=Q(u;x)+\sum_{3\le j<p}C_j(u;x)+C_p(u;x).
\end{equation}
Thus degree two, the intermediate degrees, and degree $p$ remain distinct.

For $c\in F$, Lean proves the elementary-symmetric homogeneity identity
\begin{equation}\label{eq:lean-wildodd-esymm-homogeneous}
 s_j\bigl((c\text{ viewed in }K)x\bigr)=c^j s_j(x).
\end{equation}
It follows exactly that
\begin{equation}\label{eq:lean-wildodd-quadratic-homogeneous}
 Q\bigl(u;(c\text{ viewed in }K)x\bigr)=c^2Q(u;x),
\end{equation}
so the surviving term is genuinely homogeneous quadratic.

The representative-indexed wrapper
$\texttt{wildOddUpperCorrectionAtRepresentative}$ records that
\eqref{eq:lean-wildodd-correction} contains no stationary representative.
For every representative $\beta\in K$ and every translation $c_0\in K$,
\begin{equation}\label{eq:lean-wildodd-representative-translation}
 C_{\beta+c_0}(u;x)=C_\beta(u;x)=C(u;x).
\end{equation}
This is stronger than the restricted translations used later by
\texttt{ResidualCoefficients}.  It concerns only the norm-correction
factor and does not discard the simultaneously translated elementary
stationary factor or the affine translation of the critical function.

The upper depth is derived from exact conductor data rather than assumed.
Let
\[
 \chi_K=\chi_F\circ N_{K/F},\qquad
 m_F(\chi_F)=2d+1,\qquad m_K(\chi_K)=2d_K+1,
\]
and assume that $\chi_F$ has minimal conductor in its norm-character
orbit.  The exact norm-pullback formulas below the break, at the minimal
critical endpoint, and strictly above the break give, respectively,
\begin{equation}\label{eq:lean-wildodd-conductor-depth-cases}
 2d_K=
 \begin{cases}
  2d,&2d<t,\\
  t,&2d=t,\\
  t+p(2d-t),&t<2d.
 \end{cases}
\end{equation}
In the high row Lean uses the explicit identity with the full break factor
\[
 m_K=p\,m_F-(p-1)(t+1)=t+p(2d-t)+1.
\]
The critical row uses the minimal-orbit theorem, not an unstated
noncancellation assumption.

For every $3\le j<p$, the numerical lemma proves the exact two inequalities
\begin{equation}\label{eq:lean-wildodd-two-numerical-inequalities}
 T\le j(d_K+a),\qquad T\le jd_K+pa.
\end{equation}
The $p=3$ case is the empty range above.  For $p\ge5$, Lean treats
$2d<t$, $2d=t$, and $t<2d$ separately; it does not merge a case boundary
or replace $T$ by a smaller target depth.

Now let $u\in K^\times$ satisfy $v_K(u)=a$ and let
$x\in\mathfrak p_K^{d_K}$.  The exact trace-lattice input has
$D=(p-1)T$, and the wild symmetric-function theorem gives
\begin{align*}
 v_F(s_j(ux))
 &\ge \left\lfloor
   \frac{j(d_K+a)+(p-1)T}{p}\right\rfloor,\\
 v_F\bigl(N_{K/F}(u)s_j(x)\bigr)
 &\ge a+\left\lfloor
   \frac{jd_K+(p-1)T}{p}\right\rfloor.
\end{align*}
Combining the exact floors with
\eqref{eq:lean-wildodd-two-numerical-inequalities} proves, for precisely
$3\le j<p$,
\begin{equation}\label{eq:lean-wildodd-higher-membership}
 s_j(ux)\in\mathfrak p_F^T,
 \qquad N_{K/F}(u)s_j(x)\in\mathfrak p_F^T.
\end{equation}
Together with the exact endpoint cancellation and
\eqref{eq:lean-wildodd-correction-decomposition}, this yields the exact
critical-depth congruence
\begin{equation}\label{eq:lean-wildodd-quadratic-reduction}
 C(u;x)\equiv Q(u;x)\pmod{\mathfrak p_F^T}.
\end{equation}

The structure \texttt{WildOddHigherSymmetricVanishingResult} packages the
two separate memberships in \eqref{eq:lean-wildodd-higher-membership}, the
degree-$p$ equality, the exact correction reduction, the zero degree-$0$
and degree-$1$ coefficients, quadratic homogeneity, and representative
translation invariance.  The principal theorem
\texttt{wildOdd\_higherSymmetric\_vanish} derives the residue-characteristic
identity and exact trace-lattice input from the positive lower break and
the chosen generating uniformizer, derives
\eqref{eq:lean-wildodd-conductor-depth-cases} from the pullback-conductor
formulas, and constructs this package.

\LeanFile{LanglandsFirstMainLemma/Cases/WildOdd/HigherSymmetricVanishing.lean}
\lean{LanglandsFirstMainLemma.wildOddUpperDepth}
\lean{LanglandsFirstMainLemma.wildOddUpperShift}
\lean{LanglandsFirstMainLemma.wildOddUpperCorrectionTerm}
\lean{LanglandsFirstMainLemma.wildOddUpperCorrection}
\lean{LanglandsFirstMainLemma.wildOddUpperQuadraticCorrection}
\lean{LanglandsFirstMainLemma.wildOddUpperCorrectionCoefficient}
\lean{LanglandsFirstMainLemma.wildOddUpperCorrectionCoefficient_zero}
\lean{LanglandsFirstMainLemma.wildOddUpperCorrectionCoefficient_one}
\lean{LanglandsFirstMainLemma.wildOddUpperCorrectionCoefficient_two}
\lean{LanglandsFirstMainLemma.wildOddUpperCorrectionTerm_degree}
\lean{LanglandsFirstMainLemma.wildOdd_intermediate_range_empty_degree_three}
\lean{LanglandsFirstMainLemma.elementarySymmetric_algebraMap_mul}
\lean{LanglandsFirstMainLemma.wildOddUpperQuadraticCorrection_homogeneous}
\lean{LanglandsFirstMainLemma.wildOddUpperCorrection_decompose}
\lean{LanglandsFirstMainLemma.wildOddUpperCorrectionAtRepresentative}
\lean{LanglandsFirstMainLemma.wildOddUpperCorrectionAtRepresentative_translate}
\lean{LanglandsFirstMainLemma.WildOddUpperDepthFormula}
\lean{LanglandsFirstMainLemma.WildOddUpperDepthFormula.eq_piecewise}
\lean{LanglandsFirstMainLemma.wildOdd_upperDepth_eq_of_conductors}
\lean{LanglandsFirstMainLemma.wildOdd_upper_numerical_inequalities}
\lean{LanglandsFirstMainLemma.wildOdd_intermediate_terms_mem}
\lean{LanglandsFirstMainLemma.WildOddHigherSymmetricVanishingResult}
\lean{LanglandsFirstMainLemma.wildOdd_higherSymmetric_vanish}
\leanok
\SourceText{Lemma lem:odd-upper-higher-vanishing.}
\uses{mod:ramification-symmetricbounds,mod:ramification-pullbackconductors,mod:localfield-extension}
\FormalizationContract
The intermediate range is exactly $3\le j<p$; the endpoint $j=p$ is an
exact norm cancellation, not merely a term known to lie in the conductor
lattice.  The target lattice is exactly $\mathfrak p_F^{t+1}$.  The proof
uses the two displayed floor bounds with the full different contribution
$(p-1)(t+1)$ and derives the three upper-depth rows from the exact
norm-pullback conductor formulas, including the critical minimal-orbit
case and the high break factor.

Degrees $0$, $1$, $2$, $3\le j<p$, and $p$ are represented separately.
The correction coefficient is literally zero in degrees $0$ and $1$, its
surviving term is homogeneous of degree $2$, and no affine or constant term
is introduced during the reduction.  The correction is independent of the
stationary representative, so the result survives every later permitted
representative translation; this does not authorize discarding the
simultaneously translated elementary factor or critical affine phase.
Every norm multiplier is definitionally $N_{K/F}(u)$, and every positivity,
residue-degree, generating-uniformizer, conductor, valuation, and
minimal-orbit hypothesis needed by the exact imported formulas is explicit.
\end{lemma}
